% Pursuing Stacks (cotes 134-2 to 134-8): archive cote 134-4, site page 36 (PDF page 37), typescript page 124.
% typescript. Draft compared: ps3.tex:1792-1856.
% AI-transcribed by Claude Opus 5.5 (claude-opus-5-5, adaptive thinking, effort high), Message Batches API, batch ps2 (msgbatch_01NVwr1Yuo8qJf5rykeGBXPK), 2026-11-12; not yet checked by eye.
\scanpage{134-4/36}
\mspage{124}

\uline{Remark s}. The crucial assumptions here seem to be b) (which allows \deleted{to \xout{defin}}\inserted{definition}
\deleted{define}\inserted{of} a $\hand{\pi}_0$ functor, and topological intuition tied up with connectedness
to enter into play\hand{)}, and e), which implies that with respect to cartesi-%
an products, the usual intuitive background for connectedness, rooted in
the example of M = (topological spaces), is indeed valid. This condition
is clearly stronger than d), which is a mere condition for convenience
in itself (otherwise, a decomposition of e into connected components
would mean a corresponding decomposition of M as a product\inserted{category}, and every -%
thing could be looked at ``componentwise''\hand{)}. As for \unsure{c|e}), it could probably
be dispensed with, by \xout{interpreting}\inserted[typed]{still defining} $\hand{\pi}_0$(X) as a strict pro-set. In the
case of modelizers anyhow, such generalization seems quite besides the
point.

The condition e), however innocent-looking in terms of topological
intuition, seems to me an extremely strong condition indeed. I suspect
that in case M is a topos, it is equivalent to total asphericity. In any
case, if M is the topos associated to a \inserted[typed]{locally connected}topological space, we've seen
time ago \xout{this is} that \xout{the} the condition \xout{implies} e) implies X is irredu-%
cible, and hence totally aspheric. In view of this exactingness of e),
I'll not use it unless specifically stated.

3.4.

52. Before pursuing the review of ``pure homotopy notions'' begun in yes-%
terday's notes, I would like first to correct some inaccuracies which
\over{\illegible{1}}{flew} in when l\over{o}{o}oking at the relationship between the two first basic
homotopy notions, namely the notion of a homotopy relation, and the \xout{xxx}
notion of a homotopism structure. \inserted{As usual,} \deleted{The inaccuracies were caused by the}
\deleted{fact that while starting on the notes,}\inserted{\deleted{came from}} \hand{the} provisional image I had in mind
\deleted{being}\inserted{was} still somewhat vague, \hand{while} the reasonable expectations came out more
clearly through the process of writing things down (including factual in-%
accuracies !).

The two notions clearly correspond to two kind of ways of construc-%
ting new categories M' in terms of a given one, and a functor
\[ M \longrightarrow M' \]
which is bijective o\over{\unsure{f}}{n} objects, and has the property moreover that for
any category C, the corresponding functor
\[ \underline{Hom}(M',C) \longrightarrow \underline{Hom}(M,C) \]
is a fully faithful embedding in the strict sense, namely injective on

\margin{36}
